\begin{frame}{분할: ``격자 위치마다 붙은 분류기''}
  10.2의 분류 네트워크를 끝까지 펼쳐보자 --- VGG 계열처럼
  224$\times$224 입력이 2$\times$2 최대풀링을 \textbf{5번} 거치면
  공간 크기는 $224/2^5 = 7$ 로 줄고, 마지막 합성곱 층의 출력은
  7$\times$7 \kb{격자}(채널 512개). \textbf{분류 헤드}란 그
  격자를 \textbf{펼쳐서} 1,000개 클래스 점수 하나로 모으는 층일
  뿐이다.
  \vskip0.5em
  \begin{block}{의미론적 분할 헤드}
    1개의 헤드 대신, \textbf{49개의 격자 위치마다 ``클래스 1,000개''
    출력을 \textbf{공유 가중치}로 붙인다} --- 마지막 층을
    ``512$\to$1,000 선형층''이 아니라 \textbf{출력 채널 1,000개
    의 1$\times$1 합성곱}으로 바꾸면, 격자 \textbf{각 위치}에서
    512개 특징을 1,000개 클래스 점수로 변환하고 그 가중치는
    49개 위치가 \textbf{공유}(1$\times$1 합성곱이 ``위치당
    분류기'' --- 10.1 FAQ의 ``채널만 선형 결합''). 파라미터는
    $512\times1{,}000+1{,}000$ \textbf{하나분}만 추가되지만,
    답의 형태가 ``벡터 1개''에서 \textbf{7$\times$7 클래스 지도}로
    바뀜.
  \end{block}
  {\footnotesize ``각 픽셀''이 아니라 7$\times$7 격자점인 이유:
  풀링이 공간 정보를 반으로 줄인 5번의 결과 --- 실제 분할
  모델은 풀링을 줄이고 U-Net처럼 해상도를 \textbf{되살리는}
  구조를 쓴다.}
\end{frame}
